% FILE: 00_project_identity.tex
% PROJECT: crosswalks
% ROLE: type-law-crosswalk-between-standards-and-execution

\section{Project Identity}
\label{sec:crosswalks-identity}

\subsection{Purpose}
\label{sec:crosswalks-identity-purpose}

The \texttt{crosswalks/} directory is the translation layer of the process intelligence research
foundry. Its purpose is to define, with formal precision, how structural information is preserved
or destroyed when a process model or event log is projected from one formalism to another. Each
crosswalk document specifies: (1) whether the projection is lossless or lossy, (2) if lossy, which
information is destroyed and under what named \texttt{LossPolicy}, and (3) the Rust type-law
contract that enforces this policy in \texttt{wasm4pm-compat} at compile time and at admission time.

The eight crosswalk documents collectively address every major translation surface in the process
intelligence stack: \textsc{BPMN} to WF-net, \textsc{OCEL} 2.0 to \textsc{XES}, Process Tree to
WF-net, WF-net to \textsc{POWL}, Declare to \textsc{WASM} witness automata, \textsc{PM4Py} Python
types to \texttt{wasm4pm-compat} Rust structs, \textsc{ALIVE\_001} gate criteria to
\textsc{ISO/IEC} compliance standards, and \textsc{BPMN} structural components to Declare
conformance constraints via \textsc{LTLf} lattice algebra.

Without this translation layer, a process intelligence system can silently discard object types,
structural constraints, and causal links while appearing to produce valid output. The crosswalks
make that destruction visible, named, and refusable.

\subsection{Repository Surface}
\label{sec:crosswalks-identity-surface}

The crosswalks project occupies the directory \texttt{/Users/sac/process-intelligence/crosswalks/}
within the process-intelligence research foundry. It contains exactly eight Markdown documents,
each defining one crosswalk mapping. The directory was assessed on 2026-05-31 as part of the
\textsc{ALIVE\_001} gate evaluation. Criterion 9 of the \texttt{ALIVE\_GATE\_ASSESSMENT.md}
checkpoint set a target of $\geq 3$ files; the actual count of 8 was verified as \textbf{MET}.

The \texttt{project\_manifest.yaml} records the following detected surfaces: \texttt{type-law},
\texttt{conformance-checking}, and \texttt{standards-compliance}. The manifest also records that
no \texttt{README.md}, no receipt file, and no ontology file are present directly in the
\texttt{crosswalks/} directory at the time of assessment.

The \texttt{ggen-unified-emitted-artifact-census.md} assigns a \textsc{CROSSWALK\_RECEIPT} to
the directory and states that \textsc{GGEN\_ECOSYSTEM\_INTEL\_ALIVE\_001} gate 2 validates
crosswalk completeness. The receipt is described as ``implicit; per-crosswalk dating'' --- no
explicit \textsc{BLAKE3}-signed receipt chain exists for this directory as a unit.

\subsection{Primary Research Role}
\label{sec:crosswalks-identity-role}

The thesis role assigned to this project is \textbf{type-law-crosswalk-between-standards-and-execution}
(source: \texttt{project\_manifest.yaml}, field \texttt{likely\_thesis\_role}). This role
positions the crosswalks project as the formal interface between the heterogeneous world of
process mining standards --- \textsc{BPMN} 2.0, \textsc{OCEL} 2.0, \textsc{XES} 1849,
Declare, WF-nets, \textsc{POWL}, Process Trees --- and the type-safe, memory-safe execution
surface of the \texttt{wasm4pm} WebAssembly engine.

The crosswalks are not implementation code. They are research specifications of projection laws.
They define what \emph{must} be enforced at the \texttt{wasm4pm-compat} API boundary: which
projections are lawful, which are lossy by necessity, which must be refused when loss is not
permitted, and how loss is quantified when it is permitted. The downstream authorization to
implement these laws in \texttt{wasm4pm-compat} flows from this research layer.

\subsection{Key Artifacts}
\label{sec:crosswalks-identity-artifacts}

The eight source files that constitute the crosswalks project are:

\begin{enumerate}
  \item \texttt{BPMN\_TO\_WFNET.md} --- Lossy projection with \texttt{BpmnToWfNetRefusal::InclusiveGatewayNotRepresentable}.
  \item \texttt{OCEL\_TO\_XES.md} --- Inherently lossy projection; \texttt{XesExportRefusal} named enum; \texttt{ProjectionName} newtype audit trail.
  \item \texttt{PROCESS\_TREE\_TO\_WFNET.md} --- Lossless projection; \texttt{TypedLoopNode<ARITY>} compile-time Leemans law enforcement.
  \item \texttt{WFNET\_TO\_POWL.md} --- Conditionally lossless; \texttt{WfNetToPoWlRefusal::NonBlockStructured} on non-block-structured nets.
  \item \texttt{declare\_to\_wasm\_mapping.md} --- Declare template to \textsc{LTLf} to Rust witness struct mapping; lattice join (\(\sqcup\)) with \texttt{LatticeViolation}.
  \item \texttt{pm4py\_to\_wasm4pm\_data\_types.md} --- \textsc{PM4Py} Python class to \texttt{wasm4pm-compat} Rust newtype crosswalk; referential integrity invariants.
  \item \texttt{alive\_001\_to\_iso\_compliance.md} --- \textsc{ALIVE\_001} gates to \textsc{ISO/IEC} 23745, 9001:2015, 27001, and 12207:2017.
  \item \texttt{bpmn\_to\_declare\_conformance.md} --- \textsc{BPMN} gateway structure to Declare constraint via \textsc{LTLf} lattice algebra.
\end{enumerate}
